% concatenated from 2026-06-15-qbrem-arxiv.tex
 %\documentclass[10pt,a4paper]{article}
\documentclass[reqno]{amsart}

\usepackage{amssymb}
\usepackage{enumitem}
%\usepackage{ntheorem} % statt amsthm
%\usepackage{fancyhdr}    % fancyheadings ist obsolete
%\usepackage[sumlimits,intlimits]{amsmath}
% links Nummern, Summen und Integrale i
\usepackage{hyperref}
\usepackage{xcolor}


\setlength\parindent{0pt}

\sloppy

%\setlist[enumerate]{wide,labelindent=0pt}
%label=\arabic*.

\newtheorem{thm}{Theorem}[section]
\newtheorem*{thmn}{Theorem}
\newtheorem*{lemn}{Lemma}
\newtheorem*{exn}{Example}
\newtheorem*{corn}{Corollary}
\newtheorem*{defno}{Definition}
\newtheorem{cor}[thm]{Corollary}
\newtheorem{lem}[thm]{Lemma}
\newtheorem{prop}[thm]{Proposition}
\theoremstyle{definition}
\newtheorem{defn}[thm]{Definition}
\newtheorem{defnr}[thm]{Definition \& Remark}
\newtheorem{ex}[thm]{Example}
\newtheorem{rem}[thm]{Remark}
\newtheorem{quest}[thm]{Question}

\numberwithin{equation}{section}

\newcommand{\TP}[1]{\textcolor{red}{[#1]}}

\newcommand*{\Chi}{\protect\raisebox{2.5pt}{$\chi$}}

\newcommand{\nbb}{\mathbb{N}}
\newcommand{\cbb}{\mathbb{C}}
\newcommand{\rbb}{\mathbb{R}}
\newcommand{\qbb}{\mathbb{Q}}
\newcommand{\lbb}{\mathbb{L}}
\newcommand{\acal}{\mathcal{A}}
\newcommand{\ecal}{\mathcal{E}}
\newcommand{\ccal}{\mathcal{C}}
\newcommand{\ocal}{\mathcal{O}}
\newcommand{\ucal}{\mathcal{U}}
\newcommand{\vcal}{\mathcal{V}}
\newcommand{\wcal}{\mathcal{W}}
\newcommand{\mcal}{\mathcal{M}}
\newcommand{\ncal}{\mathcal{N}}
\newcommand{\lcal}{\mathcal{L}}
\newcommand{\bcal}{\mathcal{B}}
\newcommand{\hcal}{\mathcal{H}}
\newcommand{\fcal}{\mathcal{F}}
\newcommand{\pcal}{\mathcal{P}}
\newcommand{\USC}{\mathcal{USC}}
\newcommand{\qbrem}{$q$-Bremermann}

\newcommand{\ph}{pluri\-harmonic}
\newcommand{\sh}{sub\-harmonic}
\newcommand{\subh}{sub\-harmonic}
\newcommand{\sph}{sub\-pluri\-harmonic}
\newcommand{\psh}{pluri\-sub\-harmonic}
\newcommand{\pshy}{pluri\-sub\-harmonicity}
\newcommand{\spsh}{strictly pluri\-sub\-harmonic}
\newcommand{\pssh}{pluri\-super\-harmonic}
\newcommand{\spssh}{strictly pluri\-super\-harmonic}
\newcommand{\qpsh}{$q$-pluri\-sub\-harmonic}
\newcommand{\qpshy}{$q$-pluri\-sub\-harmonicity}
\newcommand{\wqpsh}{weakly $q$-pluri\-sub\-harmonic}
\newcommand{\sqpsh}{strictly $q$-pluri\-sub\-harmonic}
\newcommand{\pshh}[1]{$#1$-plurisubharmonic}
\newcommand{\PSH}{\mathcal{PSH}}

\newcommand{\psc}{pseudo\-convex}
\newcommand{\spsc}{strictly pseudo\-convex}
\newcommand{\qpsc}{$q$-pseudo\-convex}
\newcommand{\pscc}[1]{$#1$-pseudoconvex}
\newcommand{\qspsc}{strictly $q$-pseudo\-convex}

\newcommand{\pscy}{pseudo\-con\-vexity}
\newcommand{\qpscy}{$q$-pseudo\-con\-vexity}
\newcommand{\spscy}{strict pseudoconvexity}

\newcommand{\cl}[1]{\overline{#1}}
\newcommand{\nbh}{neighbourhood}

\newcommand{\cont}{continuous}
\newcommand{\usc}{upper semi-continuous}
\newcommand{\lsc}{lower semi-continuous}
\newcommand{\fct}{function}
\newcommand{\fcts}{functions}

\renewcommand{\and}{\quad \mathrm{and} \quad}
\newcommand{\bnd}{b}
\newcommand{\dd}{\partial}
\newcommand{\dbar}{\overline{\partial}}
\newcommand{\cld}{\overline{D}}
\newcommand{\dq}{\bar{\partial}}
\newcommand{\relc}{\Subset}
\newcommand{\x}{\times}
\newcommand{\lnorm}{\left|\left|}
\newcommand{\rnorm}{\right|\right|}

\newcommand{\D}{\displaystyle}
\newcommand{\eps}{\varepsilon}
\newcommand{\vphi}{\varphi}
\newcommand{\vrho}{\varrho}
\renewcommand{\Leftrightarrow}{\Longleftrightarrow}   

\newcommand{\qand}{\quad\mathrm{and}\quad}

\textheight=8.23in
%\renewcommand{\baselinestretch}{1.3}

\title[Perron-Bremermann envelope for $q$-plurisubharmonic functions]{The Perron-Bremermann envelope for $q$-plurisubharmonic functions on unbounded domains in $\cbb^n$}

\author{Pawlaschyk, Thomas}
%\address{School of Mathematics and Natural Sciences, University of Wuppertal, Gau{\ss}str. 20, 42119 Wuppertal, Germany}
%\email{pawlaschyk@uni-wuppertal.de}


%\thanks{}

\subjclass{32F10, 32U05}

\keywords{Dirichlet problem, Dirichlet-Bremermann problem, plurisubharmonic, $q$-plurisubharmonic, $q$-convex, unbounded domains, Perron-Bremermann envelope}

%\date{\today}

\begin{document}
\begin{abstract}

Let $D$ be an unbounded domain in $\mathbb{C}^n$, and let $f$ be a bounded continuous function prescribed on the boundary of $D$. We show that if $D$ has $r$-peak points on its boundary and is of bounded type, $f$ extends to a maximal bounded continuous function $F$ on $\overline{D}$ that is $q$-plurisubharmonic and $(n-q-1)$-plurisuperharmonic (i.e., $(dd^c F)^n=0$) on $D$, and that coincides with the Perron-Bremermann envelope created with respect to bounded $q$-plurisubharmonic functions on $\overline{D}$.


\end{abstract}

\maketitle

%\section*{An anecdote involving Kang-Tae Kim}

%Schockschwerenot!\footnote{An old-fashioned German exclamation of surprise, comparable to ''Um Himmels Willen!{}`` or ''Potzblitz!{}``.} My first research stay in South Korea as a Ph.D. student was in 2014 at POSTECH University in Pohang. I went there together with Martin Sera, who was also a Ph.D. student at the time, and it was there that we met Kang-Tae Kim for the first time. Even though Kang-Tae Kim was busy leading the GAIA Center, he was a perfect host — not only showing enthusiastic interest in our freshly started research, but also introducing us to Korea and its culture. It was a magnificent trip and the starting point of a developing friendship between Pohang and Wuppertal. After returning to Germany, I graduated in 2015 and immediately applied for a DAAD–NRF Scientific Exchange Program to return to POSTECH in 2016. I stayed in One of my talks was scheduled for the evening seminar starting at 8 p.m. on September 12. I stayed in my former colleague Tobias Harz's office while he was a postdoc there. I remember how I was preparing my talk on the subject of the given paper when we were suddenly shocked by a loud noise. What was it? We looked into the hallway and saw more and more heads peeking out of offices. What was happening was the foreshock of the 2016 Gyeongju earthquake, measuring 5.1 on the local magnitude scale at 7:44 p.m. Korean time. Apart from a few cracks in the walls, nothing serious happened. So it was decided that I would give my talk as planned. I had just presented a couple of classical results when the actual earthquake struck, measuring 5.8 on the local magnitude scale. Some people took cover under the tables, others stood frozen and stared at the ceiling. Kang-Tae Kim remained calm and replied with a smile, “Oh, this result was ground shaking!” We all had a good laugh. Thank you for your hospitality, for sharing your wisdom, your entertaining stories, and your friendship. Danke sch\"on!

\section*{Acknowledgements}

I developed the main ideas of this paper during my research stay in 2016 at POSTECH in Pohang, South Korea. I would like to thank Kang-Tae Kim for his hospitality during this stay and his friendship developed from many more stays over the years. I would also like to thank my doctoral supervisor and colleague Nikolay Shcherbina for having supported my academical career. He passed away on February 7, 2026, after a long, severe illness. 


\section{Introduction}\label{intro}

The Dirichlet-Bremermann problem for \psh\ \fcts\ is of special interest in pluripotential theory and several complex variables. It was solved by Bremermann~\cite{Br} for bounded \spsc\ domains with continuous boundary data using the Perron envelope. Continuity of the solution is due to Walsh~\cite{Wa}. Bedford and Taylor~\cite{BT} showed that it is the unique solution satisfying $(dd^c F)^n=0$. Recent developments on bounded domains can be found in~\cite{Nil} and~\cite{NTT}. For unbounded domains, additional assumptions are required. Simioniuc and Tomassini~\cite{SiTo} presented a constructive method for obtaining a continuous solution on strictly convex and  special strictly pseudoconvex domains. Nguyen and Hung~\cite{NH} proved results on unbounded $B$-regular domains and domains of so-called \textit{bounded type} using Jensen measures.

In this note, we study the Dirichlet-Bremermann problem for \textit{\qpsh}\ \fcts\ in the sense of~\cite{HM} from the viewpoint of the \textit{$q$-Perron-Bremermann envelope} given by
\[
\pcal_{f,q,D,M}(z)=\sup\{\psi(z) : \psi \in \PSH_q(D) \cap{\ccal(\cld)}, \ \psi \leq f \ \text{on}\ \partial D, \ \psi\leq M \ \text{on}\ \cld\},
\]
where $D$ is a possibly unbounded domain in $\cbb^n$, $f$ is a is continuous \fct\ on $\partial D$, and $M$ is a constant with $\sup_{\partial D} f \leq M \leq +\infty$. The aim is to find a maximal \cont\ \qpsh\ extension of $f$ on $\cld$.

Hunt-Murray \cite{HM} solved the issue on bounded strictly \qpsc\ domains, extending \cont\ boundary data to a so-called \textit{$q$-Bremermann} \fct\ on $\cld$, i.e., a \fct\ that is \cont\ on $\cld$, and \qpsh\ and $(n-q-1)$-plurisuperharmonic on $D$. Kalka \cite{Ka} presented a more general version on strictly $r$-pseudoconvex domains. S\l{}odkowski~\cite{Sl84} developed approximation techniques for \qpsh\ \fcts\ and deduced the uniqueness of the solution. Simioniuc and Tomassini~\cite{SiTo} extended their results from \psh\ \fcts\ ($q=0$) to the \qpsh\ setting on unbounded strictly convex domains.

In the first section, we establish basic properties of $q$-Bremermann and $q$-maximal functions. In the second one, we study the $q$-Perron-Bremermann envelope introduced above. In the third and final section, we present classical results on the $q$-Dirichlet-Bremermann problem and our main Theorem~\ref{thm-q-Dirichlet-unbounded}. We prove that $D$ admits a unique maximal continuous extension of $f$ to $\cld$ that is $q$-Bremermann on $D$, provided that $D$ has $r$-peak points ($0 \leq r \leq q \leq n-r-1$) on its boundary and is of bounded type.

\section{Properties of q-Bremermann functions}\label{sec-prop-q-brem}

We begin by listing main properties of \qbrem\ \fcts. Let $A$ be a subset of $\cbb^n$. We write $\USC(A)$ for the set of \usc\ \fcts\ $u:A\to[-\infty,+\infty)$, and $\ccal(A)$ for the set of real-valued \cont\ \fcts\ on $A$. Unless stated otherwise, $D$ denotes a domain in $\cbb^n$, possibly unbounded. $B_r(p)$ is used for the ball centered at $p \in \cbb^n$ with radius $r>0$. By $U \relc A$ we mean that $\cl{U}$ is a compact subset of $A$.

\begin{defn} Let $f:A\to\rbb$ be a real-valued \fct\ on an arbitrary subset $A$ of~$\cbb^n$. By $f^*(z):=\limsup_{\zeta \to z} f(\zeta)$ we mean the \textit{upper semi-continuous regularization} of $f$ on $A$, and by $f_*(z):=\liminf_{\zeta \to z} f(\zeta)$ the \textit{lower semi-continuous regularization} of $f$ on $A$. We denote by $\USC(A)$ the set of all \usc\ \fcts\ on $A$, and by $\ccal(A)$ the set of all \cont\ \fcts\ on $A$.
\end{defn}

\begin{rem} $f^*$ is \usc\ and $f_*$ is \lsc, satisfying $f_* \leq f \leq f^*$ on $A$. Furthermore, $(-f)^*=-f_*$. The \fct\ $f$ is \usc\ on $A$ if and only if $f^*=f$. It is \lsc\ on $A$ if and only if $f_*=f$. Therefore, $f \in \ccal(A)$ if and only if $f_*=f=f^*$. 
\end{rem}

The following notion of \qpshy\ is in the sense of \cite{HM}.

\begin{defn} \label{def-qpsh}\label{defn-sph} Let $q\in\{ 0,\ldots,n-1\}$ and let $\psi:U\to[-\infty,+\infty)$ be an \usc\ \fct\ on an open set $U$ in $\cbb^n$.
\begin{enumerate}

%\item The \fct\ $\psi$ is called \emph{\sph} \emph{on $\Omega$} if for every ball $B \relc U$ and every \fct\ $h$ which is \ph\ on a \nbh\ of $\cl{B}$ with $\psi \leq h$ on $\partial B$ we already have that $\psi \leq h$ on $\cl{B}$.

%\item The \fct\ $\psi$ is \emph{\qpsh}\ on $\Omega$ if $\psi$ is \sph\ on $\pi \cap \Omega$ for every complex affine subspace $\pi$ of dimension $q+1$.

\item The \fct\ $\psi$ is \emph{\qpsh}\ on $U$ if for every complex $(q+1)$-dimensional affine subspace $\Pi$, every ball $B \relc U$ and every \ph\ \fct\ $h$ on defined on an open neighborhood of $\cl{B}$ with $\psi \leq h$ on $\partial B \cap \Pi$ we have that $u \leq h$ on $\cl{B}\cap \Pi$.\footnote{This type of functions was called \emph{pseudoconvex of order $n-q$} by O.~Fujita~\cite{Fu2}. Smooth \qpsh\ \fcts\ are exactly the \emph{weakly $(q+1)$-convex} ones in the sense of Grauert.}


%\item A \lsc\ \fct\ $\vphi$ on $U$ is called \itdef{\qpssh}\ on $\Omega$ if $-\psi$ is \qpsh\ on $\Omega$.

\item If $m \geq n$, every \usc\ \fct\ on $U$ is by convention $m$-\psh.

\item $\PSH_q(U)$ donotes the family of \qpsh\ \fcts\ on $U$.


\item For $q=0$, we have the classical \textit{\psh}\ \fcts, and the $(n-1)$-\psh\ \fcts\ are also called \textit{subpluriharmonic}. 

\end{enumerate}

\end{defn}

The next property follows directly from the definition of \qpshy.

\begin{prop}\label{prop-sup-inf-qpsh} Let $\{\psi_j\}_{j \in J}$ be a collection of \qpsh\ \fcts\ on an open subset $U$ of $\cbb^n$. If $\{\psi_j\}_j$ is locally bounded, then the \usc\ regularization $\Psi^*$ of $\Psi:=\sup_{j \in J}\{\psi_j\}$ is \qpsh\ on $U$. If $\{\psi_j\}_{j \in \nbb}$ is a decreasing sequence, then $\inf_j\{\psi_j\}=\lim_j \psi_j$ is \qpsh\ on $U$.
\end{prop}

The maximum principle holds true for \qpsh\ \fcts~\cite{HM}.

\begin{thm}[Maximum principle] \label{prop-qpsh-locmax} Let $q\in\{ 0,\ldots,n-1\}$ and $U \relc \cbb^n$. Then any \fct\ $\psi \in \PSH_q(U) \cap \USC(\cl{U})$ fulfills
\[
\max\{ \psi(z) : z \in \cl{U}\} = \max\{\psi(z) : z \in \partial U\}.
\]
\end{thm}

The next result is Theorem~5.1 in~\cite{Sl84}.

\begin{thm}\label{thm-add-qpsh} Let $\psi \in \PSH_q(U)$ and $\varphi \in \PSH_r(U)$. Then $\psi+\varphi \in \PSH_{q+r}(U)$.
\end{thm}

The following important technique is from~\cite{Dieu}, but was also already used in some proofs in~\cite{HM} and~\cite{Sl84}.

\begin{lem}[Glueing lemma]\label{lem-glueing}

Let $W \subset U$, $w \in \PSH_q(W)$ and $u \in \PSH_q(U)$ such that $\limsup_{\zeta \to z} w(\zeta) \leq u(z)$ for every $z \in \partial W \cap U$. Then
\[
\psi(z):=\begin{cases} \max\{w(z),u(z)\}, & z \in \cl{W} \\ u(z), & z \in U\setminus W \end{cases}
\]
is \qpsh\ on $U$.
\end{lem}

We have the following characterization of smooth \qpsh\ \fcts.

\begin{thm}\label{smooth-qpsh} Let $q \in \{0,\ldots,n-1\}$ and let $\psi$ be a $\ccal^2$-smooth \fct\ on an open subset $U$ in $\cbb^n$. Then $\psi \in \PSH_q(U)$ if and only if the complex Hessian $\left( \frac{\partial^2 \psi}{\partial z_k \partial \overline{z}_l}(p)\right)_{k,l=1}^n$ has at most~$q$ negative eigenvalues at every point $p$ in~$U$.
\end{thm}

\begin{defn}\label{defn-sqpsh} We say that $\psi$ is \textit{strictly \qpsh}\ on $U$ if it is $\ccal^2$-smooth and its complex Hessian has at most $q$ non-positive eigenvalues.
\end{defn}

Using approximation techniques from \cite{Sl84} and~\cite{Bu}, we obtain the subsequent characterization of \qpshy\ (see Corollary 3.5.5 in~\cite{TPTHESIS}).

\begin{lem}\label{lem-char-qpsh} Let $q \in \{0,1,\ldots,n-1\}$ and $U$ be open in $\cbb^n$. Then $\psi \in \PSH_q(U)$ if and only if for every open subset $V \subseteq U$ and every $\varphi \in \PSH_{n-q-1}(V) \cap \ccal(V)$ the sum $\psi+\varphi$ is subpluriharmonic on $V$.
\end{lem}

The next notion is from~\cite{Sl84}. 

\begin{defn}\label{defqBrem} Let $q\in\{ 0,1,\ldots,n{-}1\}$ and $A$ be an arbitrary set in $\cbb^n$. Let $F:A\to\rbb$ be real-valued. We say that $F$ is \textit{almost \qbrem} on $A$, if $F^* \in \PSH_{q}(A^\circ)$ and $-F_* \in \PSH_{n-q-1}(A^\circ)$. Here, $A^\circ$ denotes the interior of $A$. The \fct\ $F$ is called \textit{\qbrem}\ on $A$ if $F$ is \cont\ on $A$ and almost \qbrem\ on~$A^\circ$.
\end{defn}

\begin{rem} Every \qbrem\ \fct\ is almost \qbrem, and every \cont\ almost \qbrem\ \fct\ is \qbrem, since $F^*=F$ and $F_*=F$ if $F$ is \cont. If $F$ is \qbrem\ and $\ccal^2$-smooth, then its complex Hessian has at least one zero eigenvalue. In this sense, \qbrem\ \fcts\ $F$ satisfy the complex Monge-Amp\`ere equation $(\partial\cl{\partial} F)^n=0$. By this reason, \qbrem\ \fcts\ were called \textit{$q$-complex Monge-Amp\`ere} or \textit{$q$-CMA} in~\cite{HM}. Therefore, every \ph\ \fct\ on an open subset $U$ of $\cbb^n$ is $q$-Bremermann on $U$ for every $q\in\{0,1,\ldots,n-1\}$. If $F$ is $0$-Bremermann, then $F^*$ is \psh\ and $-F_*$ is subpluriharmonic.
\end{rem}

\begin{ex} Let $\|z\|$ be any complex norm on $\cbb^n$. Then $\log\|z\|$ is \psh\ on $\cbb^n$ and $-\log\|z\|$ is subpluriharmonic on $\cbb^n\setminus\{0\}$ (see \cite{TPESZ}). This means that $\log\|z\|$ is $0$-Bremermann on $\cbb^n\setminus\{0\}$. As a consequece, $\log\|F\|$ is $0$-Bremermann on $\{F \neq 0\}$ for every holomorphic map $F:U \to \cbb^k$ and $U \subset \cbb^n$ open.
\end{ex}

\begin{prop}\label{prop-sup-inf-qbrem} Let $D$ be a domain in $\cbb^n$. If $\{F_j\}_j$ is a locally bounded increasing sequence of \qbrem\ \fcts\ on $\cld$, then $F:=\sup_j F_j$ is almost \qbrem\ on $\cld$. If $\{G_j\}_j$ is a locally bounded decreasing sequence of \qbrem\ \fcts\ on $\cld$, then $G:=\inf_j G_j$ is almost \qbrem\ on~$\cld$. 
\end{prop}

\begin{proof} By Proposition~\ref{prop-sup-inf-qpsh}, $F^* \in \PSH_q(D)$. Since $F_j$'s are \cont, $F$ is \lsc\ on $D$, so $-F$ is \usc. Moreover, $-F=\inf_j\{-F_j\}$, so $-F$ is the limit of a decreasing sequence of \fcts\ $-F_j \in \PSH_{n-q-1}(D)$. Again, by Proposition~\ref{prop-sup-inf-qpsh}, $-F=-F_*$ is also in $\PSH_{n-q-1}(D)$.

Now notice that, in general, a \fct\ $H$ is (almost) \qbrem\ if and only if $-H$ is (almost) $(n-q-1)$-Bremermann. Then $\{-G_j\}_j$ is an increasing sequence of $(n-q-1)$-Bremermann \fcts. Thus, by the first part, $-G=\sup_j\{-G_j\}$ is almost $(n-q-1)$-Bremermann on $D$, so that $G$ is $q$-Bremermann on~$D$.
\end{proof}

The next result can be found in \cite{Sl84}.

\begin{prop}[Uniqueness property]\label{prop-unique-qbrem} Let $D$ be a bounded domain in $\cbb^n$, and let $F,G$ be two almost \qbrem\ \fcts\ on $\cl{D}$ such that $F,G \in \ccal(\partial D)$ and $F=G$ on $\partial D$. Then $F=G$ on $\cld$.
\end{prop}

\begin{proof} By Thoerem \ref{thm-add-qpsh} and since $F=F_*=F^*$ and $G=G_*=G^*$ on $\partial D$, $F^*-G_*$ and $G^*-F_*$ are $(n-1)$-\psh\ on $D$ such that $F^*-G_*= 0$ and $G^*-F_*=0$ on $\partial D$. Since $F=G$ on $\partial D$, we can apply the maximum principle (Theorem~\ref{thm-add-qpsh}) in order to obtain $F\leq F^* \leq G_* \leq G$ and $G\leq G^* \leq F_* \leq F$ on~$\cld$.
\end{proof}

This motivates to study the maximal property of \qbrem\ \fcts\ in the subsequent sense.

\begin{defn} Let $D$ be a domain in $\cbb^n$ and $F:\cld \to \rbb$ a real-valued \fct. We say that $F$ is \textit{$q$-maximal on $\cld$} if $F^*\in\PSH_q(D)\cap\USC(\cld)$ and if for every $\psi \in \PSH_q(D)\cap\USC(\cl{D})$ with $\psi \leq F_*$ on $\partial D$ we have that $\psi \leq F_*$ on~$\cld$. We say that $F$ is \textit{$q$-maximal on (bounded) open subsets of $U$} if it is $q$-maximal on $\cl{U}$ for every (bounded) open subset $U$ of $D$.
\end{defn}

\begin{rem} If $F$ is $q$-maximal on open subsets of $D$, then it is $q$-maximal on $\cld$ and on bounded open subsets of $D$. If $D$ is bounded, $q$-maximality on open sets and $q$-maximality on bounded open sets are the same.
\end{rem}

\begin{prop}\label{prop-equiv-glob-loc-q-max} Let $D$ be a domain in $\cbb^n$ and $F\in\ccal(\cld)$. Then $F$ is $q$-maximal on $\cld$ if and only if it is $q$-maximal on open subsets of $D$.
\end{prop}

\begin{proof} Let $F$ be $q$-maximal on $\cld$ and let $U$ be an open subset of $D$. We pick a \fct\ $\psi \in \PSH_q(U)\cap \USC(\cl{U})$ with $\psi \leq F$ on $\partial U$. Then by the Glueing Lemma~\ref{lem-glueing},
\[
\Psi := \begin{cases} \max\{\psi,F\}, & \text{on}\ \cl{U}\\ F, & \text{on}\ \cld\setminus U \end{cases}
\]
is in $\PSH_q(U)\cap \USC(\cl{U})$ and fulfills $\Psi \leq F$ on $\partial D$. By the $q$-maximality of $F$, it follows that $\Psi\leq F$ on $\cld$, so that $\psi\leq\Psi \leq F$ on $U$.
\end{proof}


\begin{prop}\label{prop-a-q-brem-implies-loc-q-max} Let $D$ be a domain in $\cbb^n$ and let $F:\cld\to\rbb$ be an almost \qbrem\ \fct\ on $\cld$. Then $F$ is $q$-maximal on bounded open subsets of~$D$.
\end{prop}

\begin{proof} Let $U$ be a bounded open set in $D$ and $\psi \in \PSH_q(U)\cap\USC(\cl{U})$ with $\psi \leq F_*$ on $\partial U$. Then $\psi-F_*$ is $(n-1)$-\psh\ due to Theorem~\ref{thm-add-qpsh}, and $\psi-F_* \leq 0$ on $\partial U$. The maximum principle yields $\psi-F_*\leq 0$ on $U$. 
\end{proof}


\begin{prop}\label{prop-equiv-loc-q-max-2} Let $D$ be a domain in $\cbb^n$ and $F\in\ccal(\cld)$. Then $F$ is \qbrem\ on $\cld$ if and only $F$ is $q$-maximal on bounded open subsets of $D$.
\end{prop}

\begin{proof} Recall that $F^*=F=F_*$ if $F$ is \cont. Then the first direction follows immediately from Proposition~\ref{prop-a-q-brem-implies-loc-q-max}.

Now assume that $F$ is $q$-maximal on $\cld$. By the definition of $q$-maximality, we already have $F \in \PSH_q(D)$, so that it only remains to show that $-F \in \PSH_{n-q-1}(D)$. Assume that this is not the case. According to Lemma~\ref{lem-char-qpsh}, there exist a ball $B \relc D$ and a \cont\ \qpsh\ \fct\ $\varphi$ defined on an open \nbh\ of $\cl{B}$ such that $-F+\varphi<0$ on $\partial B$ and $-F(p)+\varphi(p)>0$ for some $p \in B$. This means that $\varphi < F$ on $\partial B$ and $\varphi(p)>F(p)$. Now define
\[
\Phi:=\begin{cases} \max\{F,\varphi\}, & \text{on}\ \cl{B}\\ F, & \text{on}\ \cld\setminus B \end{cases}
\]
By the Glueing Lemma~\ref{lem-glueing}, the \fct\ $\Phi \in \PSH_q(D) \cap \ccal(\cld)$ satisfying $\Phi = F$ on $\partial D$ and $\Phi(p)=\varphi(p)>F(p)$. This contradicts the $q$-maximality of $F$ on $\cld$. The proof is finished due to Proposition~\ref{prop-equiv-glob-loc-q-max}.
\end{proof}






\section{Properties of the q-Perron-Bremermann envelopes}

The solution of the Dirichlet-Bremermann problem for \psh\ \fcts\ uses the Perron-Bremermann envelope which easily extends to \qpsh\ \fcts. Notice that, in this section, $M \in \rbb\cup\{+\infty\}$, unless otherwise stated.

\begin{defn}\label{defn-q-Perron}  Let $D$ be a domain in $\cbb^n$, and $f \in \ccal(\partial D)$.  Fix a constant $M$ with $\sup_{\partial D} f \leq M \leq +\infty$. Then the \textit{$q$-Perron-Bremermann envelope} on $\cld$ is defined at each point $z$ of $\cld$ by
\[
\pcal_{f,q,D,M}(z):=\sup\{\psi(z) : \ \psi \in \PSH_q(D)\cap \ccal(\cld), \ \psi \leq f \ \text{on}\ \partial D, \ \psi \leq M \ \text{on}\ \cld\}.
\]
If $M=+\infty$, we simply write $\pcal_{f,q,D}$ rather than $\pcal_{f,q,D,+\infty}$ because in this case, the condition $\psi \leq M$ is superfluous.
\end{defn}

\begin{rem} It is obvious that $\pcal_{f,q,D,M} \leq M$ and $\sup_{M<+\infty} \pcal_{f,q,D,M} \leq \pcal_{f,q,D}$. Notice that $\pcal_{f,q,D,M} \leq f$ on $\partial D$, but the \usc\ regularization $\pcal_{f,q,D,M}^*$ might jump at the boundary and exceed $f$ there. This means that we cannot conclude immediately that $\pcal_{f,q,D,M}^*=\pcal_{f,q,D,M}$. If the latter identity holds true, $\pcal_{f,q,D,M}$ is \cont\ on $\cld$.

If $D$ is bounded, by the maximum principle, the assumption $\psi \leq f \leq M$ on $\partial D$ implies $\psi \leq M$ on~$\cld$. In this case, for any $M \geq \sup_{\partial D} f$, we have, $\pcal_{f,q,D,M}  = \pcal_{f,q,D} \leq \sup_{\partial D} f$. 

Since $\pcal_{f,q,D,M}$ is \lsc, it is $q$-maximal with respect to $\psi \in \PSH_q(D) \cap \ccal(\cld)$ such that $\psi \leq f$ on $\partial D$ and $\psi \leq M$ on $\cld$.
\end{rem}

\begin{rem}\label{rem-slightly} In the literature, the definition of the $q$-Perron-Bremermann envelope ($q=0$) is slightly different. E.g., in \cite{NH} or \cite{NTT} (for $q=0$) the following is used:
\[
\widetilde{\pcal}_{f,q,D}(z):=\sup\{\psi(z) : \ \psi \in \PSH(D), \ \psi^{\star} \leq f \ \text{on}\ \partial D\},
\]
where $\psi^{\star}(z):=\limsup_{\zeta \to z}\psi(\zeta)$ for $z \in \cld$. It is clear that it coincides with the definition in \cite{HM},
\[
\sup\{\psi(z) : \ \psi \in \PSH(D)\cap\USC(\cld), \ \psi \leq f \ \text{on}\ \partial D\}.
\]
The reason why we chose our definition is, that in contrary to $\pcal_{f,q,D}$, the \fct\ $\widetilde{\pcal}_{f,q,D}$ need not to be \lsc\ on $D$, which we use in several arguments. In general, we have $\pcal_{f,q,D} \leq \widetilde{\pcal}_{f,q,D}$. We will discuss conditions on their equality later.
\end{rem}


\begin{prop}\label{prop-q-envelope} Let $D$ be a domain in $\cbb^n$. The $q$-Perron-Bremermann envelope $\pcal:=\pcal_{f,q,D,M}$ is \lsc\ with $\pcal \leq f$ on $\partial D$, and it is almost \qbrem\ outside $E_{\infty}:=\{z \in \cld : \pcal(z)=+\infty\}$ in $\cld$. Moreover, it is \cont\ and \qbrem\ outside $E$ in $\cld$, where $E:=E_\infty \cup E_d$ and $E_d:=\{z \in \cld : \pcal(z) < \pcal^*(z)\}$. Hence, $\pcal$ is $q$-maximal on bounded open subsets of $D\setminus E$. Since $f \in \ccal(\partial D)$, we have $E \cap \partial D = \emptyset$. If $M<+\infty$ or if $D$ is bounded, then $E_\infty=\emptyset$.
\end{prop}

\begin{proof} $\pcal$ is the supremum of \cont\ \fcts\ on $\cld$, so it is \lsc, i.e., $\pcal_*=\pcal$, and $\pcal^* \in \PSH_q((D\setminus E_\infty)^\circ)$. It remains to show that $-\pcal$ is $(n-q-1)$-\psh\ on $(D\setminus E_\infty)^\circ$. Assume that this is not the case. Then by Lemma~\ref{lem-char-qpsh} there exist a ball $B \relc D\setminus E_\infty$ and a \fct\ $\varphi \in \PSH_q(B) \cap \ccal(\cl{B})$ such that $-\pcal + \varphi < 0$ on $\partial B$, but $-\pcal(p) + \varphi(p) >0$ for some point $p \in B$. In other words, $\vphi<\pcal\leq M$ on $\partial B$ and $\pcal(p) < \varphi(p)$. By the maximum principle, $\varphi(p)\leq M$. By the continuity of $\varphi$, the lower semi-continuity of $\pcal$ and the definition of the $q$-Perron-Bremermann envelope, there exists a $\psi \in \PSH_q(D) \cap \ccal(\cld)$ such that $\pcal \geq \psi > \varphi$ on $\partial B$,  $\psi(p)\leq\pcal(p) < \varphi(p)$ and $\psi \leq f$ on $\partial D$. Now define
\[
\Psi:= \begin{cases} \max\{\psi, \varphi\}, & \text{on}\ \cl{B}\\ \psi, & \text{on}\ \cld\setminus B \end{cases}.
\]
By the Glueing Lemma~\ref{lem-glueing}, $\Psi \in \PSH_q(D) \cap \ccal(\cld)$, $\Psi=\psi \leq f$ on $\partial D$, $\Psi(p)= \varphi(p) > \pcal(p)$ and $\Psi \leq M$ on~$\cld$. By the definition of the $q$-Perron-Bremermann envelope, we obtain $\Psi \leq \pcal$ on $\cld$. In particular, $\Psi(p) \leq \pcal(p)$, which is a contradiction. Hence, $-\pcal$ is $(n-q-1)$-\psh\ on $(D\setminus E_\infty)^\circ$. The rest is obvious, since $\pcal=\pcal_*$ and $E_d$ is the set of discontinuity of $\pcal$ on $\cld$.
\end{proof}

\begin{ex} Let $n \in \nbb$, $z \in \cbb$, and define
\[
s_n(z):=n(|\exp(\exp(z))|-1)=n(e^{e^x\cos(y)}-1),
\]
on $D:=\rbb\times(-\pi/2,\pi/2) \subset \cbb$. Then $s_n=0$ on $\partial D$ and $s_n>0$ on $D$. Let $f:=0$ on $\partial D$. Then, for every $n \in \nbb$, we have $s_n \leq \pcal_{f,q,D}$. Therefore, $\pcal_{f,q,D}=+\infty$ on $D$, and $E_\infty=D$. By the Riemann mapping theorem and the maximum principle, $\pcal_{f,q,D,M}=0$ on $\cld$ for any $M \geq 0$. Hence, $\sup_{M<+\infty}\pcal_{f,q,D,M} < \pcal_{f,q,D}$ on $D$. 
\end{ex}


\begin{prop}\label{prop-q-Perron-q-max} Let $D$ be a domain in $\cbb^n$.

\begin{enumerate}

\item If there exists a \fct\ $F$ on $\cld$ such that $F \in \ccal(\cld)$ and $F$ is $q$-maximal with respect to $\psi \in \PSH_q(D)\cap\USC(\cld)$ with $\psi \leq f$ on $\partial D$ and $\psi \leq M$ on~$\cld$, where $f:=F|_{\partial D}$ and $M\geq\sup_{\partial D} f$, then $\pcal_{f,q,D,M} = F$ is \cont\ on $\cld$.

\item $\pcal_{f,q,D,M} \in \ccal(\cld)$ is $q$-maximal on bounded open subsets of $D$. 
\end{enumerate}

\end{prop}

\begin{proof} (1) Obviously, $F\leq \pcal_{f,q,D,M}$ on $\cld$. On the other hand, for every $\psi$ as above in (1), by $q$-maximality of $F$, we have $\psi \leq F$ on $\cld$. Thus, $\pcal_{f,q,D,M} \leq F$ on~$\cld$, and therefore, $\pcal_{f,q,D,M} = F$ on $\cld$.

(2) Since $\pcal_{f,q,D,M} \in \ccal(\cld)$, it is \qbrem\ on $\cld$, and therefore $q$-maximal on bounded open subsets of $D$ due to Proposition~\ref{prop-equiv-loc-q-max-2}.
\end{proof}







\section{Solution of the q-Dirichlet-Bremermann problem}

The $q$-Perron-Bremermann envelope attains the prescribed values at the boundary when the domain permits peak points at its boundary.

\begin{defn}\label{defn-q-peak} An open set $U$ in $\cbb^n$ has a \textit{$q$-peak point} at a boundary point $p \in \partial U$ if there is a \textit{$q$-peak \fct}\ $\psi \in \PSH_q(U) \cap \ccal(\cl{U})$ such that $\psi(p)=0$ and $\psi<0$ on $\cl{U}\setminus\{p\}$.
\end{defn}

\begin{rem} Every \textit{local} $q$-peak point for $U$ is a $q$-peak point for $U$. Indeed, assume that there is an open \nbh\ $V$ of $p$ such that $U \cap V$ has a $q$-peak point at $p$. Take an open ball $B \relc V$ centered at $p$ and let $C:=\max_{\partial B} \psi$. Then $\varphi:=\max\{\psi,a\}$ extends by the constant $C<a<0$ from $\cl{U} \cap \cl{B}$ to a \qpsh\ \fct\ on the whole of $U$ that is \usc\ on $\cl{U}$ and peaks at~$p$.
\end{rem}

The next lemma clarifies the relevance of $q$-peak points for the $q$-Perron-Bremermann envelope.

\begin{lem}\label{lem-q-peak} Let $p \in \partial D$ be a $q$-peak point for the domain $D$ in $\cbb^n$ with $q$-peak function $\psi$, and let $f \in \ccal(\partial D)$ with $f \geq0$. Then $\pcal_{f,q,D,M}(p) = f(p)$ for every $M \geq \sup_{\partial D} f$.
\end{lem}

\begin{proof} For every $\eps>0$ there is a radius $r>0$ such that $\Psi_{a,\eps}:=a\psi+f(p)-\eps < f$ on the closure of the ball $B=B_r(p)$ in $\partial D$ for all $a\geq1$. If $a\geq1$ is large enough, $\Psi_{a,\eps} < 0$ on $\partial B \cap \cld$. By the Glueing Lemma~\ref{lem-glueing}, $\widetilde{\Psi}_{a,\eps}:=\max\{0,\Psi_{a,\eps}\} \in \PSH_q(D)\cap\ccal(\cld)$ such that $\widetilde{\Psi}_{a,\eps} \leq f$ on $\partial D$ and $\widetilde{\Psi}_{a,\eps} \leq \sup_{\partial D} f \leq M$ on $\cld$. Therefore, $\widetilde{\Psi}_{a,\eps}(p)=f(p)-\eps \leq \pcal_{f,q,D,M}(p) \leq f(p)$. Since $\eps>0$ was chosen arbitrarily, we conclude $\pcal_{f,q,D,M}(p) = f(p)$.
\end{proof}

Domains that admit peak points at its boundary are given by the following important class of sets.

\begin{defn}\label{defn-sqpsc} We say that a domain $D$ in $\cbb^n$ is \textit{strictly (Levi) \qpsc}\ at $p \in \partial D$ if there is a $\ccal^2$-smooth \sqpsh\ \fct\ $\psi$ on an open \nbh\ $U$ of $p$ such that $D \cap U=\{z \in U : \psi<0\}$ and $d\psi(p)\neq 0$.
\end{defn}

\begin{rem}\label{rem-q-Dirichlet} The classical \spscy\ is strict $0$-pseudoconvexity in our sense. This means that the result in Theorem \ref{thm-q-Dirichlet} always holds true in the case of $r=0$ for any $q \in \{0,1,\ldots,n-1\}$. If $D$ is strictly Levi \qpsc\ at $p \in \partial D$, then it has a (local) $q$-peak point at $p$. In view of Definition~\ref{defn-sqpsc} we simply can take $\varphi:=\psi-\eps|z-p|^2$ as a local peak function for $D$ at $p$, for $\eps>0$ small enough.
\end{rem}

The following theorem is due to \cite{HM} in the case $r=q$ and in its general version to~\cite{Ka}. Uniqueness of the solution has been shown in~\cite{Sl84}. Notice also the example and the comment before Theorem 3.7 in~\cite{Bu} on the extra condition $r \leq q \leq n-r-1$. In~\cite{Dieu}, it was pointed out that $r$-peak points (or $B_q$-regular points) at the boundary of the given domain  are sufficient to get a solution of the $q$-Dirichlet-Bremermann problem.

\begin{thm}\label{thm-q-Dirichlet} Let $r \in \{0,1,\ldots,n-1\}$ and let $D$ be a bounded domain in $\cbb^n$ that admits an $r$-peak point at each of its boundary points. Then for $r \leq q \leq n-r-1$ and $f \in \ccal(\partial D)$, the $q$-Perron-Bremermann envelope $\widetilde{\pcal}_{f,q,D}$ is \cont\ on $\cld$. Moreover, $\widetilde{\pcal}_{f,q,D}$ is the unique $q$-Bremermann \fct\ on $D$ such that $\D \widetilde{\pcal}_{f,q,D} = f$ on $\partial D$. Therefore, it is $q$-maximal on open subsets of $D$.
\end{thm}



%\begin{cor} The $q$-Perron-Bremermann envelope $\pcal_{f,q,D}$ from Theorem~\ref{thm-q-Dirichlet} is $q$-maximal.\end{cor}


\begin{rem} It follows from the definition that $\pcal_{f,q,D} \leq \widetilde{\pcal}_{f,q,D}$. Since $\widetilde{\pcal}_{f,q,D}$ is \cont\ on $\cld$ and $\widetilde{\pcal}_{f,q,D}=f$ on $\partial D$, we also obtain $\widetilde{\pcal}_{f,q,D} \leq \pcal_{f,q,D}$. Therefore, $\pcal_{f,q,D} = \widetilde{\pcal}_{f,q,D}$ under the assumptions made in Theorem~\ref{thm-q-Dirichlet}.
\end{rem}



\begin{cor}\label{prop-equiv-loc-q-max} If $D$ is a domain in $\cbb^n$ and $F \in \ccal(\cld)$. $F$ is \qbrem\ on $\cld$ if and only if $F=\pcal_{F|_{\partial B},q,B}$ for every ball $B \subseteq D$. 
\end{cor}

\begin{proof}
One direction follows from the uniqueness of \qbrem\ \fcts\ (Proposition~\ref{prop-unique-qbrem}). If $F=\pcal_{F|_{\partial B},q,B}$ on $B$ for every ball $B$ in $D$, then $F$ is locally \qbrem. Since \qpshy\ is a local property, the proof is finished.
\end{proof}


\begin{rem} The boundedness of $D$ guarantees that $\pcal_{f,q,D} \leq \max_{\partial D} f < +\infty$ on $\cld$ by the maximum principle, and Walsh's technique \cite{Wa} can then be applied to establish the continuity of $\pcal_{f,q,D}$ on $\cld$. For an unbounded domain $D$, however, it is not immediate that $\pcal_{f,q,D}$ is bounded, and even when boundedness does hold, continuity remains the most delicate issue. As we shall see later, this issue plays a central role in our main result, Theorem~\ref{thm-q-Dirichlet-unbounded}. 
\end{rem}

Consider also the following Example 8.2 from~\cite{ShTo} on the boundedness of the boundary data.

\begin{ex}\label{ex-Sh-To} Consider the domain $D=\{(z,u+iv) \in \cbb^2 : v > |z|^2 + u^2\}$. It is strictly convex in $\cbb^2$, and there exists a \fct\ $f \in \ccal(\partial D)$ with $f \leq 0$ such that for each $\psi \in \PSH(\cld)$ with $\psi \leq f$ on $\partial D$ we have $\psi\equiv-\infty$ on $D$. Therefore, $\pcal_{f,{q=0},D}=-\infty$ on $D$.
\end{ex}

In virtue of the previous example, from now on, we only consider continuous boundary values that are bounded below. The next result is Theorem~24 and Theorem~25 in~\cite{SiTo}.


\begin{thm}\label{thm-SiTo} Let $D$ be an unbounded strictly convex domain in $\cbb^n$ and $f \in \ccal(\partial D)$ with $f\geq0$. Then there exists a \fct\ $\Phi \in \PSH_q(D) \cap \USC(\cld)$ such that $\Phi \geq0$, $\Phi|_{\partial D}=f$ and $\Phi \in \ccal(\partial D)$. Moreover, $\Phi$ is $q$-maximal on $\cld$. If $f$ is bounded, $\Phi$ is \cont\ and \qbrem\ on $\cld$.
\end{thm}

\begin{rem} It follows from Proposition~\ref{prop-q-Perron-q-max} that $\Phi=\pcal_{f,q,D}$ on $\cld$, and from Proposition~\ref{prop-q-envelope} that $\Phi$ is almost $q$-Bremermann on $\cld$.
\end{rem}



Now we construct a \cont\ \qpsh\ extension out of given \cont\ boundary data (see also \cite{HST}).

\begin{lem}\label{lem-qpsh-extension} Let $r \in \{0,1,\ldots,n-1\}$ and let $D$ be a domain in $\cbb^n$ that admits an $r$-peak point at each of its boundary points. Fix $r \leq q \leq n-r-1$. Then, for every $f\in\ccal(\partial D)$ with $f \geq 0$, there exists $F \in \PSH_q(D) \cap \ccal(\cld)$ such that $F|_{\partial D}=f$, $F\geq0$ and $\sup_{\cld} F = \sup_{\partial D} f$.
\end{lem}

\begin{proof} Let $\{B_i\}_i$, $\{B_i'\}_i$ and $\{B_i''\}_i$ be locally finite coverings of $\partial D$ by balls centered in $p_i \in \partial D$ with $B_i \relc B_i' \relc B_i''$. Let $D_i:=B''_i \cap D$. For each $i$, we pick a \cont\ function $f_i$ on $\overline{D_i}$ such that $f_i=f$ on $\partial D\cap\overline{B}_i$,  $f_i < 0$ on $\partial D_i \setminus \overline{B_i'}$, $f_i\leq f$ on $\partial D \cap B_i''$. Then $\sup_{\overline{D_i}} f_i = \sup_{\overline{D_i}} f$. Since all boundary points of $D_i$ are $r$-peak points, it follows from  Theorem~\ref{thm-q-Dirichlet} that for each $f_i$ there is a \qbrem\ \fct\ $F_i$ on $\overline{D_i}$ such that $F_i=f_i$ on $\partial D_i$ and $\sup_{\overline{D_i}} F_i = \sup_{\overline{D_i}} f$. We set $F_i':=\max\{ F_i, 0\}$. Then $F_i'$ is \qpsh\ on $D_i$ and \cont\ on $\cl{D_i}$, and $F_i'<0$ near $\partial D_i\setminus B_i'$. Therefore, the function $F_i'':\cld \to \rbb$ given by $F_i'':=F_i'$ on $\overline{D_i}$ and $F_i'':=0$ on $\cld\setminus B_i''$ is \cont\ on $\cld$ and \qpsh\ on $D$. Since $\{B_i\}_i$ is a locally finite covering of $\partial D$, the function $F:=\sup_i\{F_i''\}=\max_i\{ F_i''\}$ is \qpsh\ on $D$ such that $F=f_i=f$ on $\partial D \cap B_i$. It is evident that
\[
\sup_{\cld} F \leq \sup_i \sup_{\overline{D_i}} F_i'' \leq \sup_{i}\sup_{\overline{D_i}} F_i \leq \sup_{\cl{D}} f.
\]
\end{proof}

\begin{rem} Observe that $F=0$ on $D\setminus \bigcup_i B_i''$. Since the covering $\{B_i''\}_i$ can be chosen arbitrarily small, we can construct $F$ to vanish on any closed subset $A \subset D$ with $A \cap \partial D = \emptyset$.
\end{rem}

%\begin{rem} The solution of the previous proposition is not necessarily globally $q$-maximal. But if it were globally $q$-maximal, then it is necessary that $\pcal_{f,D} \leq F$. The left hand side function is known as the upper envolope of \qpsh\ \fcts\ with boundary data $f$. So if $F$ is a globally $q$ maximal \qpsh\ \fct, then $F=\pcal_{f,D}$. By this reason, in the general case, $\pcal_{f,D}$ is a natural candidate for a globally $q$-maximal \qpsh\ \fct\ on $\cl{D}$. \TP{Check arguments ??? Also locally $q$-maximal??}\end{rem}

The next notion can be found in \cite{NH}.

\begin{defn}\label{defn-bounded-type} An unbounded domain $D$ in $\cbb^n$ is of \textit{bounded type} if there exist $\Psi \in \PSH(D)\cap \ccal(D)$ such that $\Psi <0$ on $D$ and $\Psi(z)\to-\infty$ as $|z|\to+\infty$, i.e., for every $C\geq0$ there exists $R>0$ such that $\Psi < -C$ on $D\setminus B_R(0)$.
\end{defn}

\begin{ex} If a domain $D$ contains a copy of $\cbb$, then it is not of bounded type by the Liouville property for \psh\ \fcts. For example, the subsequent domain $D$ is \spsc, but not of bounded type, since it contains  $\cbb\times\{0\}$,
\[
\{(z,w) \in \cbb^2 : \log|w| + |z|^2+|w|^2<1\}.
\]
If a domain $D$ is contained in $\{z \in \cbb^n : |p(z)|^2>(1+|z|^2)^{\deg p}\}$ for a complex polynomial~$p \in \cbb[z_1,\ldots,z_n]$, then $D$ is of bounded type (see \cite{NH}). 
\end{ex}

The maximum principle holds true on unbounded domains of bounded type.

\begin{prop} \label{prop-qpsh-max-princ-2} Let $q\in\{ 0,\ldots,n-1\}$ and let $D$ be a domain in $\cbb^n$ of bounded type. Then any bounded above \fct\ $\vphi \in \PSH_q(D) \cap \USC(\cl{D})$ satisfies
\[
\sup\{ \vphi(z) : z \in \cl{D}\} = \sup\{\vphi(z) : z \in \partial D\}.
\]
\end{prop}

\begin{proof} Assume that the assumption is false, and that there exists $\vphi \in \PSH_q(D) \cap \USC(\cl{D})$ such that $
m:=\sup_{\partial D}\vphi < M:=\sup_{\cld} \vphi < +\infty$. Let $\Psi$ be from Definition~\ref{defn-bounded-type}, and $\widehat{\Psi}$ its \cont\ extension to $\cld$. We fix $\eps>0$ and define $\Phi_{\eps}:=\varphi+\eps\widehat{\Psi}$. Let $R_0>0$ be so large that $\eps\widehat{\Psi}< m-M$ on $\cld\setminus B_{R_0}(0)$ for every $R\geq R_0$. Then $\Phi_{\eps} \leq m$ on $\cld\setminus B_{R}(0)$, and by the maxmum principle on bounded domains, $\Phi_{\eps} \leq m$ on $\cld \cap \cl{B_{R}(0)}$ for every $R \geq R_0$. Therefore, $\Phi_{\eps} \leq m$ on $\cld$. Since $\eps>0$ was arbitrarily chosen, $\varphi \leq m$ on $\cld$. But this contradicts to $M\leq\sup_{\cld} \vphi \leq m$.
\end{proof}

\begin{rem}\label{rem-qpsh-max-princ-2} If $D$ is of bounded type and $f \in\ccal(\partial D)$ is bounded above, then $\pcal_{f,q,D,M} =\pcal_{f,q,D,M_0}$ for every $M$ with $M_0:=\sup_{\partial D} f \leq M < +\infty$.
\end{rem}

We present our main result on the solution of the $q$-Dirichlet-Bremermann problem on unbounded domains of bounded type.

\begin{thm}\label{thm-q-Dirichlet-unbounded}
Let $r$, $q$, $D$ be as in Theorem~\ref{thm-q-Dirichlet}, except that $D$ is unbounded. Let $f \in \ccal(\partial D)$ with $0 \leq f \leq \sup_{\partial D} f \leq M$ for some upper bound $M < +\infty$. Then there exists an almost \qbrem\ \fct\ $\Phi=\Phi_M$ on $\cld$ such that

\begin{enumerate}

\item $\Phi=f$ on $\partial D$, and $0 \leq \Phi \leq M$ on $\cld$.

\item $\Phi$ is $q$-maximal with respect to \qpsh\ \fcts\ $\psi$ with $\psi \leq f$ on $\partial D$ and $\psi \leq M$ on $\cld$. Therefore, $\pcal_{f,q,D,M} \leq \Phi$.

\item $\Phi$ is \cont\ on $\partial D$.

\item If $D$ is of bounded type, $\Phi$ is \cont\ and \qbrem\ on $\cld$. Thus, $\Phi=\pcal_{f,q,D,M_0}$ for $M_0:=\sup_{\partial D} f$.

\item If $D$ is of bounded type, $\Phi$ is the unique \qbrem\ \fct\ on $\cld$ with $\Phi=f$ on $\partial D$.% and $0\leq \Phi \leq \sup_{\partial D} f$ on $\cld$.

\end{enumerate}

\end{thm}

\begin{proof} 

\textit{Claim 1:} Existence of an \usc\ almost \qbrem\ $\Phi$ with $\Phi=f$ on $\partial D$ and $0 \leq \Phi \leq M$.

\medskip

For $j \in \nbb_0$ we define $D_j:=D\cap B_{j}(0)$,  $E_j:=\partial B_{j}(0)\cap \cld$ and $A_j:=\partial D \cap \cl{B_{j}(0)}$. We may assume that $j$ is so large that $D_0$ is not empty.

For $j \geq 1$, let $\vphi_{j} \in \ccal(\cl{D_j})$ such that $ 0 \leq f \leq \vphi_{j}\leq M$ on $\partial D_j \cap \partial D$, $\vphi_{j}=f$ on $\partial D_{j-1}\cap \partial D$ and $\vphi_j=M$ on $E_j$. Let $\Phi_{j}$ be the solution of the $q$-Dirichlet-Bremermann problem with boundary data $\vphi_{j}$ on $D_j$. Then by Theorem~\ref{thm-q-Dirichlet}, $\Phi_j=\pcal_{\vphi_j,q,D_j}$ is \qbrem\ on $\cl{D_j}$, $\Phi_j=\vphi_j$ on $\partial D_j$, $q$-maximal on open subsets of $D_j$, and fulfills $0 \leq \Phi_j \leq M$ on $\cl{D_j}$. 

Since $\Phi_{j+1}=f$ on $\partial D_j\cap \partial D$ and $\Phi_{j+1} \leq M$ on $D_j$, we have that $\Phi_{j+1} \leq \vphi_{j}$ on $\partial D_j$. Since $\Phi_{j+1}$ is \qpsh\ on $D_j$ and $\Phi_{j}$ is $q$-maximal on $\cl{D_j}$, we have that $\Phi_{j+1} \leq \Phi_{j}$ on $\cl{D_j}$. Hence, for each $j_0$, we obtain a locally bounded, decreasing sequence of \qbrem\ \fcts\ $\{\Phi_j\}_{j\geq j_0}$ on $\cl{D_{j_0}}$ such that $\Phi_j=f$ on $\partial D \cap \partial D_{j-1}$. Since $\{D_j\}_{j}$ is an increasing sequence of domains such that $D=\bigcup_{j}D_j$, the \fct\ $\Phi:=\inf_j \Phi_j$ is well-defined on the whole of $\cld$. Furthermore, it fulfills $\Phi=f$ on $\partial D$ and $0\leq\Phi \leq M$ on $\cl{D_j}$ for each $j$, so $\Phi \leq M$ on $\cld$. By Proposition~\ref{prop-sup-inf-qpsh}, $\Phi$ is almost \qbrem\ on $\cld$, and by the continuity of $\Phi_j$'s, $\Phi \in \USC(\cld)$.

\medskip

\textit{Claim 2:} $\Phi$ is $q$-maximal with respect to \qpsh\ \fcts\ $\psi$ with $\psi \leq f$ on $\partial D$ and $\psi \leq M$ on $\cld$, and $\pcal_{f,q,D,M} \leq \Phi$.

\medskip

Let $\psi \in \PSH_q(D)\cap\USC(\cld)$ such that $\psi \leq f$ on $\partial D$ and $\psi \leq M$ on $D$. Define $\psi_j:=\psi|_{\cl{D_j}}$. Then, by the construction of $\vphi_j$ and $\Phi_j$, we have $\psi_j\leq \vphi_j=\Phi_j$ on $\partial D_j$ for each $j$. Since $\Phi_j$ is $q$-Bremermann on $\cl{D_j}$, we conclude that $\psi=\psi_j \leq \Phi_j$ on $\cl{D_j}$ for every $j$, so that $\psi \leq \Phi$ on $\cl{D}$. Therefore, $\pcal_{f,q,D,M} \leq \Phi$ on~$\cld$.



\medskip

\textit{Claim 3:} $\Phi$ is \cont\ on $\partial D$.

\medskip

We only need to show that $\Phi$ is \lsc\ on $\partial D$, since it is already \usc\ on $\cld$. By Lemma~\ref{lem-qpsh-extension}, there exists a \fct\ $F \in \PSH_q(D)\cap\ccal(\cld)$ such that $F=f$ on $\partial D$ and $0 \leq F \leq \sup_{\cld} F = \sup_{\partial D} f \leq M$. By $q$-maximality of $\Phi$, we have $F\leq\Phi$ on $\cld$. Thus, if $p \in \partial D$, then
\[
\liminf_{z \to p} \Phi(z) \geq \liminf_{z \to p} F(z) = \lim_{z \to p} F(z) = F(p)=f(p)=\Phi(p)
\]

%\textit{Claim 3:} $\Phi$ is \lsc\ on $D$, i.e. for every $\eps>0$, $p \in D$ there is a $\delta>0$ such that for every $z \in B_\delta(p)$ we have $\Phi(p)-\eps < \Phi(z)$.

%\medskip

%We show for $z=p-w$ and $|w| < \delta$ we have
%\[\Phi(p)-\eps = \Phi(z+w)-\eps \leq \Psi(z) \leq \Phi(z),\]
%for some suitable \qpsh\ \fct\ on $D$ which extends $f$ continuously into $D$ and which is bounded ($0 \leq \Psi \leq M$). By maximilaty of $\Phi$ we get the last inequaltiy $\Psi \leq \Phi$.

\medskip

\textit{Claim 4:} If $D$ is of bounded type, $\Phi$ is \cont\ and \qbrem\ on $\cl D$, and $\Phi \leq\pcal_{f,q,D,M}$.

\medskip

We use similar arguments as in the proof of Lemma~2 in~\cite{SiTo}. It suffices to show that $\Phi$ is \lsc\ on $D$. Since $D$ is of bounded type, there is a \cont\ negative \psh\ function $\Psi$ on $D$ such that $\Psi(z) \to -\infty$ as $|z|\to+\infty$.

We fix a point $p_0 \in D$, an upper bound $M \in [\sup_{\partial D} f,+\infty)$ and an arbitrary $\eps>0$. We choose a constant $a>0$ so small that $-\eps/2 < a\Psi(p_0)$. We replace $\Psi$ by $a\Psi$, so that we can assume $-\eps/2 < \Psi(p_0)$. Since $\Psi$ is \cont\ on $D$, we obtain that $-\eps/2 < \Psi$ on a ball $B_{\rho}(p_0) \relc D$ for some $\rho>0$.

Let $r_2>r_1>0$ be so large that $B_{\rho}(p_0) \relc B_{r_1}(0)$ and $\Psi +M < 0$ on $\cld \setminus B_{r_1}(0)$. We set $B':=B_{r_1}(0)$, $B'':=B_{r_2}(0)$ and $D'':=D\setminus \cl{B''}$.
%\TP{We can convolve $R$ with an appropriate Dirac-\fct\ in order to smoothen $R$ on the set of all points $z \in D \cap B_{r_2}(0)$ with distance to the boundary of $D \cap B_{r_2}(0)$ at most $\delta/3$. Hence, we can assume that $R$ is smooth there. This implies that $-\eps/2 < R$ in a small ball $B_{\eta}(p)$. }
Since $\Phi$ is \cont\ on $\partial D$, it is uniformly \cont\ on the compact set $\partial D \cap \cl{B''}$. Hence, there is a $\delta''>0$ such that for every $p \in \partial D \cap \cl{B''}$ and $z \in \cld \cap \cl{B''}$ with $|p-z|<2\delta''$ we have
\[
|\Phi(p) - \Phi(z)| < \eps/4.
\]
% \Phi(z+w)-\eps/2 > \Phi(z+w)-\eps/2 + \Psi(z).
Let $D_1:=\{z \in D \cap B'' : d(z,\partial D)) \geq \delta/3\}$ and $D_2:=\{z \in D \cap B'' : d(z,\partial D)) < \delta/3\}$.


We fix a positive constant $\delta< \min\{\rho,r_2-r_1,\delta''\}$. For $w \in \cbb^n$ with $|w|<\delta/3$ we define
\[
\widetilde{\Phi}(z):= \begin{cases} \max\{\Phi(z), \Phi(z+w)+\Psi(z)-\eps/2\},  & z \in D_1  \\ \Phi(z) & z \in D_2 \cup (\cld \setminus B')  \end{cases}
\]
It is clear that, if $z \in D_1$, then $z+w \in \cld$. Now let $z \in D \cap B'$ such that $d(z,\partial D)=\delta/3$. Then there exists a $p \in \partial D$ such that $z \in B_\delta(p)$. Then $z+w \in B_{2\delta}(p)$, so that $|\Phi(p)-\Phi(z)|<\eps/2$ and $|\Phi(p)-\Phi(z+w)|<\eps/2$ due to uniform continuity of $\Phi$ on $\partial D$ discussed above. Therefore, since $\Psi<0$ on $D$,
\[
\Phi(z)> \Phi(z+w) - \eps/2 > \Phi(z+w) + \Psi(z) - \eps/2.
\]
Thus, $\widetilde{\Phi}$ is \qpsh\ on $D \cap B''$ by the Glueing Lemma~\ref{lem-glueing}.

If $z \in \partial B'' \cap \cld$ with $d(z,\partial D)\geq\delta/3$, then $z+w \in \cld \setminus B'$ and therefore, 
\[
\Phi(z+w)+\Psi(z)-\eps/2 \leq M + \Psi(z) < 0 \leq \Phi(z),
\]
so $\widetilde{\Phi}$ is \qpsh\ on $D$ and \usc\ on $\cld$ again by the Glueing Lemma~\ref{lem-glueing}. It is easy to see that $0 \leq \widetilde{\Phi} \leq M$ and $\widetilde{\Phi}=\Phi=f$ on $\partial D$. By the $q$-maximality of $\Phi$ on $\cld$ (Claim 2), we have $\widetilde{\Phi} \leq \Phi$ on $\cld$. Recall that we fixed a point $p_0 \in D$ such that $B_{\rho}(p_0) \relc B' \cap D$,   and $-\eps/2 < \Psi$ on $B_{\rho}(p_0)$. Hence, $d(p_0,\partial D) > \rho>\delta$. Then $z=p_0-w \in D_1$ for $|w| < \delta/3$, so that
\[
\Phi(p_0) -\eps =\Phi(z+w) -\eps < \Phi(z+w) + \Psi(z) -\eps/2 \leq \widetilde{\Phi}(z) \leq \Phi(z).  
\]
This means that $\Phi$ is \lsc\ at $p_0$. Since $p_0$ was an arbitrary point in $D$, we conclude that $\Phi$ is \lsc\ on the whole of $\cl{D}$. Since we already deduced that $\Phi$ is \usc\ on $\cld$, we finally get that $\Phi$ is \cont\ on the whole of $\cld$. Since $\Phi$ is almost \qbrem\ on $\cld$, it is \qbrem\ on $\cld$. Moreover, also by the continuity and the properties of $\Phi$, we have $\Phi \leq \pcal_{f,q,D,M}$. It follows from the maximum principle of domains of bounded type that $\Phi=\pcal_{f,q,D,M}=\pcal_{f,q,D,M_0}$ for every $M \geq M_0:=\sup_{\partial D} f$ (see Remark~\ref{rem-qpsh-max-princ-2}).


\medskip

\textit{Claim 5:} $\Phi=\pcal_{f,q,D,M_0}$ is the unique \qbrem\ \fct\ on $\cld$ that equals $f$ on~$\partial D$.% and $0 \leq F \leq M$ on $\cld$.

\medskip

Assume that another \qbrem\ $F$ exists on $\cld$ with $F=f$ on $\partial D$. It is clear that $F \leq \Phi$ on $\cld$ by the maximum principle on domains of bounded type and the definition of $\pcal_{f,q,D,M_0}$. Recall that $F$ is $q$-maximal on bounded open subsets of $D$, since it is \qbrem\ on $\cld$ (see Proposition~\ref{prop-equiv-loc-q-max-2}).

For arbitrary $\eps>0$ define $\Phi_\eps:=\Phi+\eps \widehat{\Psi}$ (where $\Psi$ is from Claim 4 and $\widehat{\Psi}$ is the \cont\ extension of $\Psi$ into the whole of $\cld$). If the ball $B=B_{r_0}(0)$ is large enough ($r_0>1/\eps$), then $\Phi_\eps \leq F$ on $\cld \setminus B$, so $\Phi_\eps \leq F$ on $\partial(D \cap B_r(0))$ for every $r \geq r_0$. By the $q$-maximality of $F$ on bounded open sets, we deduce $\Phi_\eps \leq F$ on $\cld$. Since $\eps>0$ was chosen to be arbitrary, we obtain $\Phi \leq F$ on $\cld$.
\end{proof}

\begin{rem} In view of Remark~\ref{rem-slightly}, for $M \geq \sup_{\partial D}f$ and $z \in D$, we define 
\[
\widetilde{\pcal}_{f,q,D,M}(z):=\sup\{\psi(z) : \ \psi \in \PSH_q(D)\cap \USC(\cld), \ \psi \leq f \ \text{on}\ \partial D, \ \psi \leq M \ \text{on}\ \cld\}.
\]
Then we can replace $\pcal_{f,q,D,M}$ by $\widetilde{\pcal}_{f,q,D,M}$ in Claim~2 in order to get $\widetilde{\pcal}_{f,q,D,M} \leq \Phi$. If $D$ is of bounded type, we get $\Phi=\pcal_{q,f,D,M}\leq \widetilde{\pcal}_{f,q,D,M} \leq \Phi$, so that $\pcal_{q,f,D,M} = \widetilde{\pcal}_{f,q,D,M}=\Phi$ in this case. 
\end{rem}

\begin{rem}
The case $M=+\infty$ appears to be more delicate. This case was studied in \cite{NH} and \cite{SiTo}. The discussion preceding Theorem 23 of the latter paper suggests that the \textit{unbounded \psh\ hull} $\cld^{\wedge}$ of the closure $\cld$ of the domain $D$ may be relevant. It is defined by
$\cld^{\wedge} := \bigcup_{r>0} \big(\cl{D \cap B_r(0)}\big)^{\wedge}$, where $K^{\wedge}$ denotes \psh\ hull of a compact set $K \subset \cld$. This motivates to define the \textit{$q$-hull} of $K$ in $\cld$ via
\[
K^{\wedge}_q := \{ z \in \cld : \psi(z) \leq \max_{K} \psi \ \text{for\ all}\ \psi \in \PSH_q(D)\cap\USC(\cld)\}.
\]
Then $K^\wedge_0=K^\wedge$, and $\pcal_{f,q,D,M} \leq \max_{\partial D \cap \cl{B_r(0)}} f \leq M$ holds true at least on the \textit{unbounded $q$-hull} $\cld^{\wedge}_q := \bigcup_{r>0} \big(\cl{D \cap B_r(0)}\big)^{\wedge}_q$ for every sufficiently large $r>0$ such that $D \cap B_r(0)\neq \emptyset$. More precisely, $\cld^{\wedge}_q \subseteq \cld \setminus E_\infty$ (see Proposition~\ref{prop-q-envelope}). In general, we have $\cld^{\wedge}_q \subsetneq \cld$, but it is not immediate that $\pcal_{f,q,D,M}$ is \cont\ on~$\cld^{\wedge}_q$, as pointed out in \cite{SiTo} in the case $q=0$.
\end{rem}
%------------------------------------------

\section*{Declerations}

This study was funded by the DAAD–NRF Scientific Exchange Program (2016). 

%\textbf{Financial interests.} The author has no relevant financial or non-financial interests to disclose.


\bibliographystyle{alpha}
\begin{thebibliography}{Tiefe}

%\bibitem[Bas76]{Ba}
%Basener, R.~F.:
%\newblock Nonlinear Cauchy-Riemann equations and q-pseudoconvexity.
%\newblock Duke Math. J., 43(1), 203--213 (1976). \href{https://doi.org/10.1215/S0012-7094-76-04317-9}{https://doi.org/10.1215/S0012-7094-76-04317-9}

\bibitem[BT82]{BT}
Bedford, E., Taylor, A.: A new capacity for plurisubharmonic functions, Acta Math. \textbf{149}, 1--40 (1982). \href{https://api.semanticscholar.org/CorpusID:123003061}{https://doi.org/10.1007/BF02392348}

\bibitem[Bre59]{Br} Bremermann, H.J.: On a generalized Dirichlet problem for plurisubharmonic functions and pseudo-convex domains. Characterization of \v{S}ilov boundaries. Trans. Amer. Math. Soc. \textbf{91}, 246--276 (1959). \href{https://doi.org/10.2307/1993121}{https://doi.org/10.2307/1993121}

\bibitem[Bun90]{Bu}
Bungart, L.: Piecewise smooth approximations to q-plurisubharmonic functions. Pacific J. Math., 142(2), 227--244 (1990). \href{http://projecteuclid.org/euclid.pjm/1102646343}{http://projecteuclid.org/euclid.pjm/1102646343}

\bibitem[Fuj92]{Fu2}
Fujita, O.: On the equivalence of the q-plurisubharmonic functions and the pseudoconvex functions of general order, Ann. Reports of Graduate School of Human Culture, Nara Women's  Univ. \textbf{7}, 77--81 (1992). 

\bibitem[HST21]{HST}
Harz, T., Shcherbina, N., Tomassini, G.: On defining functions and cores for unbounded domains. III, Sb. Math. \textbf{212} 859 (2021). \href{https://doi.org/10.1070/SM8898}{https://doi.org/10.1070/SM8898}


\bibitem[HM78]{HM}
Hunt, L.R., Murray, J.J.: $q$-plurisubharmonic functions and a generalized Dirichlet problem, Michigan Math. J. \textbf{25} (3), 299--316 (1978). \href{https://doi.org/10.1307/mmj/1029002112}{https://doi.org/10.1307/mmj/1029002112}

\bibitem[Kal79]{Ka}
Kalka, M.: On a conjecture of Hunt and Murray concerning $q$-plurisubharmonic functions, Proc. Amer. Math. Soc. \textbf{73}, 30--34 (1979). \href{https://doi.org/10.2307/2042875}{https://doi.org/10.2307/2042875}

\bibitem[Ngu06]{Dieu}
Nguyen, Q.D.: $q$-plurisubharmonicity and $q$-pseudoconvexity in $\cbb^n$, Publ. Mat., Vol. \textbf{50}, 349--369 (2006). \href{https://doi.org/10.5565/publmat\_50206\_05}{https://doi.org/10.5565/publmat\_50206\_05}

\bibitem[NH08]{NH}
Nguyen, Q.D.: Hung, D.H., Jensen measures and unbounded $B-$regular domains in ${\mathbf{C}}^n$, Annales de l'Institut Fourier, Vol. \textbf{58} no. 4, 1383--1406 (2008). \href{https://doi.org/10.5802/aif.2388}{https://doi.org/10.5802/aif.2388}

\bibitem[NTT26]{NTT}
Nguyen, Q.D.: Tang, V.L., Tran, D.H., Dirichlet Problem for Plurisubharmonic Functions on Bounded Quasi B-regular Domains in $\cbb^n$. J. Geom. Anal. \textbf{36}, 164 (2026). \href{https://doi.org/10.1007/s12220-026-02412-1}{https://doi.org/10.1007/s12220-026-02412-1}


\bibitem[Nil25]{Nil}
Nilsson, M.: Plurisubharmonic functions with discontinuous boundary behavior, Indiana Univ. Math. J. \textbf{74} No. 2, 539–553 (2025). \href{https://doi.org/10.1512/iumj.2025.74.60165}{https://doi.org/10.1512/iumj.2025.74.60165}

%\bibitem[NW21]{NW}
%Nilsson, M., Wikström, F.: Quasibounded plurisubharmonic functions, Internat. J. Math., 32(9) (2021). \href{https://doi.org/10.1142/S0129167X21500683}{https://doi.org/10.1142/S0129167X21500683}

\bibitem[Paw15]{TPTHESIS}
Pawlaschyk, T.: On some classes of $q$-plurisubharmonic functions and $q$-pseudoconcave sets. Dissertation zur Erlangung eines Doktorgrades, Bergische Universit\"at Wuppertal (2015). \href{https://nbn-resolving.org/urn:nbn:de:hbz:468-20151210-101726-1}{https://nbn-resolving.org/urn:nbn:de:hbz:468-20151210-101726-1}

\bibitem[PZ13]{TPESZ}
Pawlaschyk, T., Zeron, E.S.: On convex hulls and pseudoconvex domains generated by $q$-plurisubharmonic functions, part I.,J. Math. Anal. App. \textbf{408}, 394--408 (2013). \href{https://doi.org/10.1016/j.jmaa.2013.05.074}{https://doi.org/10.1016/j.jmaa.2013.05.074}

\bibitem[ST99]{ShTo}
Shcherbina, N., Tomassini, G.: The Dirichlet problem for Levi-flat graphs over unbounded domains, Internat. Math. Res. Notices. \textbf{1999} (3), 111--151 (1999). \href{https://doi.org/10.1155/S1073792899000069}{https://doi.org/10.1155/S1073792899000069}

\bibitem[ST08]{SiTo}
Simioniuc, A., Tomassini, G.: The Bremermann–Dirichlet problem for unbounded domains of $\mathbb{C}^n$, manuscripta math. \textbf{126}, 73–97 (2008). \href{https://doi.org/10.1007/s00229-008-0167-x}{https://doi.org/10.1007/s00229-008-0167-x}

\bibitem[S{\l}o84]{Sl84}
S{\l}odkowski, Z.: The Bremermann-Dirichlet problem for $q$-plurisubharmonic functions, Ann. Scuola Norm. Sup. Pisa Cl. Sci., Serie 4, Vol. \textbf{11} no. 2, 303--326 (1984). \href{https://www.numdam.org/item/ASNSP\_1984\_4\_11\_2\_303\_0}{https://www.numdam.org/item/ASNSP\_1984\_4\_11\_2\_303\_0}

\bibitem[S{\l}o86]{Sl86}
S{\l}odkowski, Z.: Local maximum property and $q$-plurisubharmonic functions in uniform algebras, J. Math. Anal. Appl. \textbf{115} (1), 105--130 (1986). \href{https://doi.org/10.1016/0022-247X(86)90027-2}{https://doi.org/10.1016/0022-247X(86)90027-2}

\bibitem[Wal69]{Wa}
Walsh, J.B.: Continuity of envelopes of plurisubharmonic functions, J. Math. Mech. \textbf{18} No. 2, 143--148, (1968). \href{https://www.jstor.org/stable/24901797}{https://www.jstor.org/stable/24901797}

\end{thebibliography}

\vspace{1cm}

Thomas Pawlaschyk\\
University of Wuppertal\\
School of Mathematics and Natural Sciences\\
Gau{\ss}str. 20, 42119 Wuppertal, Germany\\
\href{mailto:pawlaschyk@uni-wuppertal.de}{\texttt{pawlaschyk@uni-wuppertal.de}}\\
\texttt{\href{https://orcid.org/0009-0004-0494-3273}{https://orcid.org/0009-0004-0494-3273}}\\

\end{document}

